\documentclass[11pt]{article}
\usepackage{fontspec}
\setmainfont[
  Path=fonts/,
  Extension=.ttf,
  UprightFont=*-400,
  BoldFont=*-700
]{NotoSerifSC}
\setsansfont[
  Path=fonts/,
  Extension=.ttf,
  UprightFont=*-400,
  BoldFont=*-700
]{NotoSerifSC}
\XeTeXlinebreaklocale "zh"
\XeTeXlinebreakskip=0pt plus 1pt
\usepackage[a4paper,margin=1.70cm]{geometry}
\usepackage{amsmath,amssymb,mathtools,bm}
\usepackage{booktabs,longtable,array}
\usepackage{enumitem}
\setlength{\parindent}{2em}
\setlength{\parskip}{0.35em}
\setlist{nosep,leftmargin=2em}
\allowdisplaybreaks[3]
\numberwithin{equation}{section}
\pagestyle{plain}

\newcommand{\R}{\mathbb{R}}
\newcommand{\E}{\mathbb{E}}
\newcommand{\Pp}{\mathbb{P}}
\newcommand{\1}{\mathbf{1}}
\newcommand{\T}{\mathsf{T}}
\newcommand{\F}{\mathrm{F}}
\newcommand{\cN}{\mathcal{N}}
\newcommand{\cL}{\mathcal{L}}
\newcommand{\cS}{\mathcal{S}}
\newcommand{\cH}{\mathcal{H}}
\newcommand{\cA}{\mathcal{A}}
\newcommand{\cM}{\mathcal{M}}
\newcommand{\cT}{\mathcal{T}}
\newcommand{\diag}{\operatorname{diag}}
\newcommand{\tr}{\operatorname{tr}}
\newcommand{\Var}{\operatorname{Var}}
\newcommand{\Cov}{\operatorname{Cov}}
\newcommand{\softmax}{\operatorname{softmax}}
\newcommand{\KL}{D_{\mathrm{KL}}}
\newcommand{\sg}{\operatorname{sg}}
\begin{document}
    \begin{center}
      {\LARGE\bfseries 38\_DeepSeek LLM}
    \end{center}
    \vspace{0.8em}

    \section{符号与维度}
    \begin{longtable}{>{$}l<{$} >{$}l<{$} p{10cm}}
    \toprule
    \text{符号} & \text{维度} & \text{含义}\\
    \midrule
    N,D,C & \R_{>0} & 参数量、训练 token 数与计算预算。\\
        L(N,D) & \R & 验证损失缩放模型。\\
        \alpha,\beta & \R_{>0} & 参数与数据幂律指数。\\
        n_L,d,V & \mathbb N & 层数、隐藏宽度与词表大小。\\
    \bottomrule
    \end{longtable}

    所有向量默认是列向量；未特别说明时，期望同时对数据、噪声和采样时刻取值。下文只展开论文中承担方法逻辑的公式，并把所需假设写在推导发生的位置。

\section{可分离幂律与计算约束下的最优分配}

设验证损失模型
\begin{equation}
 L(N,D)=L_\infty+A N^{-\alpha}+B D^{-\beta},
 \qquad A,B,\alpha,\beta>0,
\end{equation}
其中 $N$ 是模型规模代理量、$D$ 是训练 token 数。训练计算近似
\begin{equation}
 C=\kappa ND.
\end{equation}
固定 $C$ 时 $D=C/(\kappa N)$，代入可约为一元函数
\begin{equation}
 \widetilde L(N)=L_\infty+A N^{-\alpha}
 +B\left(\frac{\kappa N}{C}\right)^\beta.
\end{equation}
一阶条件
\begin{equation}
 -\alpha A N^{-\alpha-1}
 +\beta B\left(\frac\kappa C\right)^\beta N^{\beta-1}=0.
\end{equation}
整理得
\begin{equation}
 N^{\alpha+\beta}
 =\frac{\alpha A}{\beta B}\left(\frac C\kappa\right)^\beta.
\end{equation}
因此
\begin{align}
 \boxed{N^*(C)=
 \left(\frac{\alpha A}{\beta B}\right)^{1/(\alpha+\beta)}
 \left(\frac C\kappa\right)^{\beta/(\alpha+\beta)},}\\
 \boxed{D^*(C)=
 \left(\frac{\beta B}{\alpha A}\right)^{1/(\alpha+\beta)}
 \left(\frac C\kappa\right)^{\alpha/(\alpha+\beta)}.}
\end{align}
两个最优指数相加为 $1$，保证 $N^*D^*\propto C$。

在最优点，一阶条件还给出两项的比例
\begin{equation}
 \alpha A(N^*)^{-\alpha}=\beta B(D^*)^{-\beta}.
\end{equation}
不是两项数值相等，而是乘上各自幂指数后相等。

\section{最优损失随计算的幂指数}

将 $N^*,D^*$ 代入，两项都按
\begin{equation}
 C^{-\frac{\alpha\beta}{\alpha+\beta}}
\end{equation}
衰减。因此
\begin{equation}
 L^*(C)=L_\infty+K C^{-\gamma},
 \qquad \gamma=\frac{\alpha\beta}{\alpha+\beta},
\end{equation}
常数 $K$ 由 $A,B,\alpha,\beta,\kappa$ 决定。
若任一资源的幂指数很小，它会成为调和式组合中的瓶颈。

边际计算收益
\begin{equation}
 \frac{dL^*}{dC}=-\gamma K C^{-\gamma-1}
\end{equation}
随规模快速减小；幂律外推给的是平均趋势，不保证单次训练的确定回报。

\section{对数线性拟合与不可忽略的 $L_\infty$}

若只拟合单项 $y=A x^{-\alpha}$，取对数得到
\begin{equation}
 \log y=\log A-\alpha\log x,
\end{equation}
可以线性回归。但实际有不可约项
\begin{equation}
 L(x)=L_\infty+A x^{-\alpha},
\end{equation}
则
\begin{equation}
 \log L\ne\log A-\alpha\log x.
\end{equation}
只有对 $L-L_\infty$ 取对数才线性，而 $L_\infty$ 本身未知，需要非线性联合拟合。

若误设 $\widetilde L_\infty=L_\infty+\delta$，变换量变为
\begin{equation}
 \log(L-\widetilde L_\infty)
 =\log(Ax^{-\alpha}-\delta).
\end{equation}
在大规模、幂律项较小时，对 $\delta$ 极敏感；这解释了外推指数和渐近损失之间的强相关。

\section{参数估计的不确定性传播}

令拟合参数
$\psi=(\log A,\log B,\alpha,\beta,L_\infty)$，最小二乘目标
\begin{equation}
 S(\psi)=\sum_i w_i[L_i-L(N_i,D_i;\psi)]^2.
\end{equation}
在最优点 $\widehat\psi$ 附近二阶近似
\begin{equation}
 S(\psi)\approx S(\widehat\psi)
 +\frac12(\psi-\widehat\psi)^TH(\psi-\widehat\psi).
\end{equation}
若残差近似高斯，参数协方差可估为
\begin{equation}
 \Cov(\widehat\psi)\approx\widehat\sigma^2H^{-1}.
\end{equation}
对由参数导出的最优规模 $g(\psi)=\log N^*(C)$，delta 方法给出
\begin{equation}
 \Var(g(\widehat\psi))
 \approx\nabla_\psi g^T\Cov(\widehat\psi)\nabla_\psi g.
\end{equation}
因此只报告单条最优曲线会隐藏幂指数拟合误差经外推放大的区间。

\section{模型规模代理量的选择}

稠密 Transformer 每 token 前向 FLOPs 常近似与参数量成正比，训练再乘常数，
于是 $C\approx\kappa ND$。MoE 中总参数 $N_{\rm total}$ 与每 token 激活参数
$N_{\rm active}$ 不同：
\begin{equation}
 C\approx\kappa N_{\rm active}D,
\end{equation}
而容量和存储更接近 $N_{\rm total}$。若把总参数直接代入稠密计算公式，会高估训练计算。

假设每层有 $E$ 个专家、每 token 激活 $k$ 个，专家参数总量 $N_e$，则粗略
\begin{equation}
 N_{\rm active}=N_{\rm shared}+\frac{k}{E}N_e,\qquad
 N_{\rm total}=N_{\rm shared}+N_e.
\end{equation}
路由不均衡还会让实际设备时间由最拥挤专家决定，因此 FLOPs 相同也不保证墙钟时间相同。

\section{两阶段或多阶段训练的计算核算}

若不同阶段模型规模、序列长度或激活稀疏度不同，总计算应相加
\begin{equation}
 C_{\rm total}=\sum_{j=1}^{J}\kappa_jN_{{\rm active},j}D_j.
\end{equation}
把全部 token 简单相加再乘最终模型规模只有在各阶段 $\kappa_jN_j$ 相同才成立。
这对继续预训练、长上下文阶段和多 token 预测辅助头尤其重要。

\section{自回归似然与 token 梯度}

对序列 $x_{1:T}$，因果语言模型按概率链式法则分解
\begin{equation}
 p_\theta(x_{1:T})=\prod_{t=1}^{T}p_\theta(x_t\mid x_{<t}).
\end{equation}
负对数似然
\begin{equation}
 \cL_{\rm NLL}=-\sum_{t=1}^{T}\log p_\theta(x_t\mid x_{<t}).
\end{equation}
若第 $t$ 步 logits 为 $z_t\in\R^{V}$，概率
$p_{tj}=e^{z_{tj}}/\sum_ke^{z_{tk}}$，真实 token 索引 $y_t$，则
\begin{equation}
 \ell_t=-z_{t,y_t}+\log\sum_ke^{z_{tk}},
\end{equation}
逐 logit 梯度
\begin{equation}
 \frac{\partial\ell_t}{\partial z_{tj}}
 =p_{tj}-\1\{j=y_t\}.
\end{equation}
Hessian 正是 softmax Jacobian
\begin{equation}
 \nabla_z^2\ell_t=\diag(p_t)-p_tp_t^T\succeq0.
\end{equation}
所以交叉熵关于 logits 凸，但关于深网参数并不凸。

对非 padding token 集合 $\mathcal I$，平均损失和困惑度
\begin{equation}
 \bar L=\frac1{|\mathcal I|}\sum_{t\in\mathcal I}\ell_t,
 \qquad \operatorname{PPL}=e^{\bar L}.
\end{equation}
因此 PPL 是真实 token 概率倒数的几何平均：
\begin{equation}
 \operatorname{PPL}=\left(
 \prod_{t\in\mathcal I}\frac1{p_\theta(x_t\mid x_{<t})}
 \right)^{1/|\mathcal I|}.
\end{equation}
不同 tokenizer 的 token 粒度不同，PPL 不能不经换算直接横向比较。

\section{旋转位置编码的相对位置恒等式}

对二维坐标对，定义旋转
\begin{equation}
 R(m\omega)=
 \begin{bmatrix}
 \cos(m\omega)&-\sin(m\omega)\\
 \sin(m\omega)&\cos(m\omega)
 \end{bmatrix}.
\end{equation}
位置 $m,n$ 的 query、key 分别为
$q_m=R(m\omega)q$、$k_n=R(n\omega)k$。内积
\begin{align}
 q_m^Tk_n
 &=q^TR(m\omega)^TR(n\omega)k\\
 &=q^TR((n-m)\omega)k.
\end{align}
第二步用到 $R(a)^T=R(-a)$ 和 $R(a)R(b)=R(a+b)$。
所以绝对位置旋转后的内积只依赖相对距离 $n-m$。

在 $d_h$ 维中对相邻坐标对使用频率
\begin{equation}
 \omega_j=\theta^{-2j/d_h},\qquad j=0,\ldots,d_h/2-1.
\end{equation}
低 $j$ 频率高、分辨近距离，高 $j$ 频率低、覆盖远距离。若推理长度远超训练长度，
相位 $m\omega_j$ 的外推和周期混叠会改变注意力几何；位置插值本质上是重新缩放这些相位。

\section{RMSNorm 的 Jacobian}

忽略可学习增益，RMSNorm
\begin{equation}
 y=\frac{x}{s},\qquad
 s=\sqrt{\frac1d x^Tx+\epsilon}.
\end{equation}
微分
\begin{equation}
 ds=\frac{x^Tdx}{d s},
\end{equation}
因此
\begin{align}
 dy&=\frac1s dx-\frac{x}{s^2}ds\\
 &=\left(\frac1sI-\frac{xx^T}{d s^3}\right)dx.
\end{align}
故
\begin{equation}
 J_{\rm RMS}=\frac1sI-\frac{xx^T}{d s^3}.
\end{equation}
与 LayerNorm 相比，它没有去均值投影
$I-d^{-1}\1\1^T$，因此不消除常数方向。加上逐坐标增益 $g$ 后，
Jacobain 左乘 $\diag(g)$。

谱上，对任意 $v\perp x$，$Jv=v/s$；沿 $x$ 方向
\begin{equation}
 Jx=\left(\frac1s-\frac{\|x\|^2}{d s^3}\right)x
 =\frac{\epsilon}{s^3}x.
\end{equation}
当 $\epsilon$ 很小，径向尺度方向被强烈压制，而切向方向按 $1/s$ 缩放。

\section{SwiGLU 前馈层的逐项梯度}

SwiGLU 可写
\begin{equation}
 a=xW_a,\qquad b=xW_b,\qquad
 h=\operatorname{swish}(a)\odot b,\qquad
 y=hW_o,
\end{equation}
其中
\begin{equation}
 \operatorname{swish}(a)=a\sigma(a).
\end{equation}
其导数
\begin{align}
 \frac d{da}[a\sigma(a)]
 &=\sigma(a)+a\sigma(a)(1-\sigma(a)).
\end{align}
令 $g_h=\partial\cL/\partial h$，则
\begin{align}
 g_a&=g_h\odot b\odot
 [\sigma(a)+a\sigma(a)(1-\sigma(a))],\\
 g_b&=g_h\odot\operatorname{swish}(a),\\
 \nabla_{W_a}\cL&=x^Tg_a,\qquad
 \nabla_{W_b}\cL=x^Tg_b,\\
 \nabla_x\cL&=g_aW_a^T+g_bW_b^T.
\end{align}
门分支 $b$ 同时调制激活和 $a$ 分支梯度，故两个投影不是可独立解释的普通 MLP。

\section{MoE 路由与负载}

对 token $x_i$，router logits $r_i=W_rx_i$，概率
\begin{equation}
 p_{ie}=\frac{e^{r_{ie}}}{\sum_{j=1}^{E}e^{r_{ij}}}.
\end{equation}
若 top-$k$ 专家集合 $S_i$，归一化门
\begin{equation}
 \widetilde p_{ie}=\frac{p_{ie}\1\{e\in S_i\}}
 {\sum_{j\in S_i}p_{ij}},
 \qquad
 y_i=\sum_{e\in S_i}\widetilde p_{ie}F_e(x_i).
\end{equation}
固定集合 $S_i$ 内可正常求导；集合边界的 arg-top-$k$ 不连续，几乎处处把未选专家梯度置零。

令软重要性
\begin{equation}
 P_e=\frac1N\sum_ip_{ie},
\end{equation}
硬负载
\begin{equation}
 F_e=\frac1N\sum_i\1\{e\in S_i\}/k.
\end{equation}
常见平衡项
\begin{equation}
 L_{\rm bal}=E\sum_eP_eF_e
\end{equation}
在均匀 $P_e=F_e=1/E$ 时等于 $1$。把硬 $F_e$ 停止梯度后，
router 梯度只通过 $P_e$，会降低当前过载专家的概率；它是训练启发式，不等同于严格容量约束。

若每专家容量为
\begin{equation}
 C_e=\left\lceil\frac{Nk}{E}\,c_{\rm cap}\right\rceil,
\end{equation}
超过容量的 token 必须丢弃或转送。即使平均负载均匀，局部批次方差仍可能溢出。

\section{低秩 KV 潜变量的吸收}

设 key/value 由共享潜变量
\begin{equation}
 c_t=W_cx_t\in\R^{d_c},\qquad
 k_t=W_kc_t,\qquad v_t=W_vc_t.
\end{equation}
注意力 logit
\begin{equation}
 q_s^Tk_t=q_s^TW_kc_t=(W_k^Tq_s)^Tc_t.
\end{equation}
因此解码时可把 $W_k$ 吸收到 query 投影，缓存 $c_t$ 而非完整 $k_t$。
若 value 聚合
\begin{equation}
 o_s=\sum_tp_{st}W_vc_t=W_v\left(\sum_tp_{st}c_t\right),
\end{equation}
$W_v$ 可放到加权和之后。缓存从每 token 的 $d_k+d_v$ 降为 $d_c$，
前提是位置编码等操作不破坏这些线性交换关系；对只作用在部分 key 维度的 RoPE，需把旋转部分另行缓存或重构。

\section{多 token 预测的联合似然}

若同一隐藏状态 $h_t$ 用 $K$ 个头预测 $x_{t+1},\ldots,x_{t+K}$，
\begin{equation}
 \cL_t=\sum_{k=1}^{K}\lambda_k
 [-\log p_{\theta,k}(x_{t+k}\mid h_t)].
\end{equation}
共享骨干梯度
\begin{equation}
 \nabla_{h_t}\cL_t=\sum_{k=1}^{K}\lambda_k g_{t,k}.
\end{equation}
其平方范数精确展开
\begin{equation}
 \left\|\sum_k\lambda_kg_{t,k}\right\|^2
 =\sum_k\lambda_k^2\|g_{t,k}\|^2
 +2\sum_{j<k}\lambda_j\lambda_k g_{t,j}^Tg_{t,k}.
\end{equation}
交叉内积为负时存在任务冲突；辅助头退火相当于训练后期逐渐去掉这些交叉项，
使最终目标回到下一 token 似然。

\end{document}
